\centerline{\bf \Huge Разнобойчик}

\begin{flushright}
        \scriptsize{}
\end{flushright}

\problem{1}
На доске написаны двузначные числа. Каждое число составное, но любые два - взаимно просты. Какое наибольшее количество чисел может быть написано на доске?

\problem{2}
На плоскости нарисованы а) 3; б) 10; различных прямых. Можно ли раскрасить получившиеся части плоскости в два цвета, чтобы области одного цвета не граничили по отрезку прямой?

\problem{3}
Среди любых десяти из шестидесяти ребят найдутся трое одноклассников. Докажите, что среди всех них найдутся 15 одноклассников.